% GENERATED FILE. Edit canonical lessons, then run npm run generate:books.
% canonical-source-hash: fnv1a64:82d1b38dc94c169a
% canonical-lessons: JA-C122-fukai, JA-C122-taisetsu, JA-C122-hitsuyou, JA-C122-kantan, JA-C122-fukuzatsu

\chapter{Deep, Precious, Necessary, Simple, Complicated}
\label{ch:ja-deep-precious-necessary-simple-complicated}
\hlchaptermodality{122}

\begin{chapteropening}
\textbf{By the end of this chapter:} \emph{I can say deep, precious, necessary, simple and complicated.}

Complete the last lesson of chapter 122: i can say deep, precious, necessary, simple and complicated.
\end{chapteropening}

\section[\texorpdfstring{\ja{ふかい} (fukai)}{ふかい (fukai)}]{\ja{ふかい} (fukai) — deep}

\label{lesson:JA-C122-fukai}



\begin{quote}
Before the new one: say the Japanese for thin, then the Japanese for shallow.
\end{quote}

\subsection*{You'll want to know: \ja{ふかい}}
\textbf{\ja{ふかい}} — \emph{fukai} — \textquotedblleft{}deep\textquotedblright{}.

\begin{grammarlens}[title={What you've built}]
One more word for this chapter.
\end{grammarlens}

\subsection*{Your turn}
\begin{itemize}
\raggedright
  \item \emph{Say it:} \emph{fukai}
  \item \emph{Say it:} \emph{fukai}, once more
  \item \emph{Say it:} say \emph{fukai}
  \item \emph{Recall:} say the Japanese for thin, then the Japanese for shallow, then say \emph{fukai} again
\end{itemize}

\begin{tcolorbox}[breakable,colback=teal!4,colframe=teal!35!black,title={Before you move on}]
What does \ja{ふかい} mean? (\textquotedblleft{}Deep\textquotedblright{}.) Say it once more.
\end{tcolorbox}
\section[\texorpdfstring{\ja{たいせつ} (taisetsu)}{たいせつ (taisetsu)}]{\ja{たいせつ} (taisetsu) — precious}

\label{lesson:JA-C122-taisetsu}



\begin{quote}
Before the new one: say the Japanese for shallow, then the Japanese for deep.
\end{quote}

\subsection*{You'll want to know: \ja{たいせつ}}
\textbf{\ja{たいせつ}} — \emph{taisetsu} — \textquotedblleft{}precious\textquotedblright{}.

\begin{grammarlens}[title={What you've built}]
One more word for this chapter.
\end{grammarlens}

\subsection*{Your turn}
\begin{itemize}
\raggedright
  \item \emph{Say it:} \emph{taisetsu}
  \item \emph{Say it:} \emph{taisetsu}, once more
  \item \emph{Say it:} say \emph{taisetsu}
  \item \emph{Recall:} say the Japanese for shallow, then the Japanese for deep, then say \emph{taisetsu} again
\end{itemize}

\begin{tcolorbox}[breakable,colback=teal!4,colframe=teal!35!black,title={Before you move on}]
What does \ja{たいせつ} mean? (\textquotedblleft{}Precious\textquotedblright{}.) Say it once more.
\end{tcolorbox}
\section[\texorpdfstring{\ja{ひつよう} (hitsuyou)}{ひつよう (hitsuyou)}]{\ja{ひつよう} (hitsuyou) — necessary}

\label{lesson:JA-C122-hitsuyou}



\begin{quote}
Before the new one: say the Japanese for deep, then the Japanese for precious.
\end{quote}

\subsection*{You'll want to know: \ja{ひつよう}}
\textbf{\ja{ひつよう}} — \emph{hitsuyou} — \textquotedblleft{}necessary\textquotedblright{}.

\begin{grammarlens}[title={What you've built}]
One more word for this chapter.
\end{grammarlens}

\subsection*{Your turn}
\begin{itemize}
\raggedright
  \item \emph{Say it:} \emph{hitsuyou}
  \item \emph{Say it:} \emph{hitsuyou}, once more
  \item \emph{Say it:} say \emph{hitsuyou}
  \item \emph{Recall:} say the Japanese for deep, then the Japanese for precious, then say \emph{hitsuyou} again
\end{itemize}

\begin{tcolorbox}[breakable,colback=teal!4,colframe=teal!35!black,title={Before you move on}]
What does \ja{ひつよう} mean? (\textquotedblleft{}Necessary\textquotedblright{}.) Say it once more.
\end{tcolorbox}
\section[\texorpdfstring{\ja{かんたん} (kantan)}{かんたん (kantan)}]{\ja{かんたん} (kantan) — simple}

\label{lesson:JA-C122-kantan}



\begin{quote}
Before the new one: say the Japanese for precious, then the Japanese for necessary.
\end{quote}

\subsection*{You'll want to know: \ja{かんたん}}
\textbf{\ja{かんたん}} — \emph{kantan} — \textquotedblleft{}simple\textquotedblright{}.

\begin{grammarlens}[title={What you've built}]
One more word for this chapter.
\end{grammarlens}

\subsection*{Your turn}
\begin{itemize}
\raggedright
  \item \emph{Say it:} \emph{kantan}
  \item \emph{Say it:} \emph{kantan}, once more
  \item \emph{Say it:} say \emph{kantan}
  \item \emph{Recall:} say the Japanese for precious, then the Japanese for necessary, then say \emph{kantan} again
\end{itemize}

\begin{tcolorbox}[breakable,colback=teal!4,colframe=teal!35!black,title={Before you move on}]
What does \ja{かんたん} mean? (\textquotedblleft{}Simple\textquotedblright{}.) Say it once more.
\end{tcolorbox}
\section[\texorpdfstring{\ja{ふくざつ} (fukuzatsu)}{ふくざつ (fukuzatsu)}]{\ja{ふくざつ} (fukuzatsu) — complicated}

\label{lesson:JA-C122-fukuzatsu}



\begin{quote}
Before the new one: say the Japanese for necessary, then the Japanese for simple.
\end{quote}

\subsection*{You'll want to know: \ja{ふくざつ}}
\textbf{\ja{ふくざつ}} — \emph{fukuzatsu} — \textquotedblleft{}complicated\textquotedblright{}.

\begin{grammarlens}[title={What you've built}]
One more word for this chapter.
\end{grammarlens}

\subsection*{Your turn}
\begin{itemize}
\raggedright
  \item \emph{Say it:} \emph{fukuzatsu}
  \item \emph{Say it:} \emph{fukuzatsu}, once more
  \item \emph{Say it:} say \emph{fukuzatsu}
  \item \emph{Recall:} say the Japanese for necessary, then the Japanese for simple, then say \emph{fukuzatsu} again
\end{itemize}

\begin{tcolorbox}[breakable,colback=teal!4,colframe=teal!35!black,title={Before you move on}]
What does \ja{ふくざつ} mean? (\textquotedblleft{}Complicated\textquotedblright{}.) Say it once more.
\end{tcolorbox}
